In recent years, the deployment of Unmanned Aerial Vehicles (UAVs) in high-risk scenarios and Non-Destructive Testing (NDT) has gained significant relevance in robotics research~\cite{siciliano2008springer}. Structural health monitoring tasks, such as the routine inspection of aging bridges and civil infrastructures, require reliable physical interaction for sensor positioning and structure monitoring, where aerial robots can effectively substitute human workers in hazardous and high-elevation workspaces~\cite{hao2024bilateral, dangelo2025semi_platform}. However, achieving fully autonomous execution in unstructured environments remains an open challenge due to environmental uncertainties, dynamic obstacles, and unexpected mission anomalies. 
\\To bridge the gap between human cognitive judgment and robotic execution, shared control (SC) and shared autonomy (SA) architectures have become the predominant paradigms~\cite{selvaggio2021autonomy}. Shared control allows a remote operator to take high-level decisions while the autonomous system remains stable and ensures constraint satisfaction. In bilateral teleoperation frameworks~\cite{hokayem2006bilateral}, incorporating haptic feedback provides the remote pilot with an intuitive physical rendering of environmental interaction forces and artificial safety barriers, as demonstrated by Hao et al.~\cite{hao2024bilateral} and D'Angelo et al.~\cite{dangelo2025semi}. 
\\Despite the advantages of shared control, performing aerial contact tasks is heavily constrained by perceptual limitations. Traditionally, aerial teleoperation relies on an egocentric perspective, where the operator supervises the scene exclusively through an onboard First-Person View (FPV) camera. During physical contact maneuvers against vertical surfaces, this perspective inevitably leads to total visual occlusion, complete loss of 3D depth perception and loss of situational awareness. Furthermore, standard underactuated quadrotors must tilt to move and to exert forces, causing significant video shake and induced motion sickness for the human operator. As demonstrated by ~\cite{uhde2024allocentric} in multi-robot coordination, shifting from an egocentric to an allocentric (or exocentric) perspective drastically reduces the operator's cognitive load, eliminates occlusions, and restores proper spatial understanding.
\\To provide an active exocentric perspective, the observing aerial platform must continuously replan its trajectory to keep the interacting agent inside its Field of View (FOV). Model Predictive Control (MPC) represents the established standard in the literature to solve constrained multi-objective aerial tracking problems. Works like~\cite{jacquet2021motor} demonstrated that perception-constrained Nonlinear MPC (NMPC) can explicitly enforce camera FOV limitations while simultaneously respecting physical actuation bounds. Similarly, in~\cite{penin2018vision} it has been addressed reactive target tracking while preventing visual occlusions and in~\cite{shayan2025exponential} it has been demonstrated that predictive formulations incorporating jerk-level constraints ensure reactive yet remarkably smooth quadrotor motion.
\\Building upon these established standards, the framework developed in this thesis addresses the residual limitations of single-agent inspection systems. For this reason, instead of relying on a single aerial manipulator, the mission is decoupled between a dedicated Inspection UAV executing the physical contact and an Observer UAV acting as an active exocentric camera. This completely separates the video stream from contact vibrations and dynamic tilting.
\\Moreover, the proposed NMPC actively tracks a moving UAV and the optimization cost explicitly penalizes linear/angular accelerations, jerk, and snap, guaranteeing smooth camera movements specifically tuned to prevent operator visual fatigue. The NMPC directly commands body thrust and torques to the PX4 autopilot in offboard mode.
\\Also, the architecture incorporates online Point of View (PoV) adjustments commanded by the operator via a haptic interface. Here the geometric FOV boundaries and safety distances are mapped into continuous repulsive forces via Artificial Potential Fields (APF), guiding the operator away from visual loss without overriding their commands.
\\For the inspection drone, stable surface contact is achieved through a high-level admittance filter, ensuring compliant interaction on standard underactuated platforms without requiring complex tiltable hardware.
